\section{Discussion and Limitations}\label{sec:discussion}

\subsection{What each layer buys}
Reading the matrix (\cref{tab:scenarios,tab:goalsummary}) top to bottom gives a
clear ordering. Transaction binding is the decisive step for RQ1: it turns the
recipient and amount swaps from \xmark{} to \cmark{} (S2a, S3, S4, S5), which a
plain OTP cannot do. Moving from a symmetric code to an asymmetric signature
(RQ2) adds two properties: resilience to a database leak (S10) and
third-party--verifiable evidence (S11), because the secret never leaves the
phone. The trusted display (RQ3) is independent of both: \texttt{sig} and
\texttt{sig\_blind} share identical cryptography, yet \texttt{sig\_blind} fails
S2a and S3 --- the strongest cryptography is wasted if the user is shown nothing
to check. The careless-user column (S2b) is the sobering bound: no mechanism in
this study defends a user who approves without looking.

\subsection{Limitations}
\paragraph{Attentive-user model.}
The attentive user is idealised: always reads, always notices. Real users sit
between this and the careless model, and the literature shows a meaningful share
of tampered transactions are approved, especially when only a digit or two of the
account changes~\cite{alzomai2008,haupert2019}. The \cmark{} cells of S2a, S3 and
G7 are therefore an \emph{upper bound} on security; we did not run a user study,
and we do not infer a detection rate from our own trials.

\paragraph{Guessing limit is per transaction (G4).}
As quantified in \cref{sec:eval}, the lockout is per transaction with no per-user
counter and no cap on transaction creation, so $k$ transactions give about
$k\cdot p$ cumulative success. This $k\cdot p$ figure is an \emph{idealised}
analytic bound that assumes the valid-code window stays fixed over the whole
guessing campaign; for \texttt{baseline} the TOTP code rotates every $30$~s, so
tens of thousands of transactions inside a single window is not a measured,
realistic capability --- only the per-transaction probabilities in
\cref{tab:bruteforce} are measured. This is weak for \texttt{baseline} and is
closed only by a per-user/per-device counter and a creation rate limit, which
introduce a denial-of-service trade-off.

\paragraph{Verifiable evidence is not non-repudiation (G6).}
The signature is checkable by a third party, but because the authenticator is
served by the bank and the public key is attested by no independent party, the
bank is not cryptographically bound by it. We claim only verifiable evidence.

\paragraph{Symmetric keys in the clear (G5).}
$K_{\text{totp}}$ and $K_{\text{hmac}}$ are stored in plaintext to make S10
concrete. A real deployment must hold them in an HSM; this is why \texttt{hmac}
is \xmark{} for S10/S11 here.

\paragraph{Trusted, separate phone (A6).}
The whole scheme rests on the phone being trusted and separate from the PC. This
is fragile: PC--phone cross-infection~\cite{dmitrienko2014}, cross-platform
synchronisation~\cite{konoth2016}, and running the bank app and the authenticator
on one phone~\cite{haupert2018} all erase the gap. A6 (phone malware) is out of
scope, and the authenticator being a web page served by the bank is itself a
trust assumption worth stating.

\paragraph{Enrolment and other out-of-scope items.}
Enrolment assumes a clean, trust-on-first-use environment, and re-enrolment in
the demo requires wiping the database; a real system needs an out-of-band
recovery process. Push notifications, native apps, WebAuthn/SPC, multiple devices
per user, and CSRF defences beyond \texttt{SameSite=Lax} are out of scope.

\subsection{Positioning}
None of the four mechanisms is novel (\cref{sec:background}); transaction binding
and WYSIWYS are long established and, in the EU, mandated as dynamic
linking~\cite{eu2018389}. The value of this work is the controlled comparison on
one code base and the blind-signing experiment that isolates the display.
Regulatory context for Vietnam (Decision 2345/Q\DJ-NHNN and any instrument that
supersedes it) should be cited only after checking the currently effective text.
